\documentclass[11pt,a4paper]{article}
\usepackage[a4paper,margin=2.4cm]{geometry}
\usepackage{amsmath,amssymb,mathtools}
\usepackage{fontspec}
\usepackage{polyglossia}
\setmainlanguage{russian}
\setotherlanguage{english}
\setmainfont{Liberation Serif}
\setsansfont{Liberation Sans}
\setmonofont{Liberation Mono}
\newfontfamily\cyrillicfont{Liberation Serif}
\newfontfamily\cyrillicfontsf{Liberation Sans}
\newfontfamily\cyrillicfonttt{Liberation Mono}
\usepackage{booktabs}
\usepackage{array}
\usepackage{enumitem}
\usepackage[hidelinks]{hyperref}

\setlength{\parindent}{0pt}
\setlength{\parskip}{0.55em}
\emergencystretch=1.5em
\setlist{itemsep=0.25em,topsep=0.3em}

\newcommand{\G}{\mathcal{G}_k}
\newcommand{\Shalf}{S^{1/2}}
\newcommand{\Smhalf}{S^{-1/2}}
\newcommand{\norm}[1]{\lVert #1\rVert_2}
\DeclareMathOperator*{\argmin}{arg\,min}
\DeclareMathOperator{\diag}{diag}

\title{Задача 1(c): условие нулевого шага в R-SGD}
\author{Homework Mini-Project: SGD versus Robust SGD under Gradient Uncertainty}
\date{16 ноября 2026 г.}

\begin{document}
\maketitle

\section*{Постановка и ответ}

Условие (3) выполнено тогда и только тогда, когда робастный шаг нулевой:
$d_k = 0$, то есть $\theta_{k+1} = \theta_k$. Из (3) \emph{не следует}
$\nabla f(\theta_k) = 0$: оно означает лишь, что нулевой вектор лежит в множестве
неопределённости $\G$, то есть данные не позволяют исключить нулевой градиент при
радиусе $\rho$.

Условие из задания:
\begin{equation*}
g_k^T S_k^{-1} g_k \le \rho^2. \tag{3}
\end{equation*}

Обозначения: $g = g_k$, $S = S_k$, $\alpha = \alpha_k$. Используется результат
пункта (b): шаг $d_k$ есть точка минимума функции $\varphi$,
\begin{equation*}
\varphi(d) \;=\; \max_{h\in\G}\Big\{h^Td + \frac{\norm{d}^2}{2\alpha}\Big\}
\;=\; g^Td + \rho\sqrt{d^TSd} + \frac{\norm{d}^2}{2\alpha},
\qquad d_k = \argmin_d \varphi(d).
\end{equation*}

Так как $S \succ 0$, матрицы $\Shalf$ и $\Smhalf$ существуют, симметричны и
положительно определены, и $\sqrt{d^TSd} = \norm{\Shalf d}$. Заметим также, что
$\varphi(0) = 0$.

Главный результат, который доказывается в шагах 1–3:
\begin{equation*}
\boxed{\; d_k = 0 \iff g_k^TS_k^{-1}g_k \le \rho^2 \iff 0\in\G \;}
\end{equation*}

Структура доказательства:

\begin{center}
\begin{tabular}{@{}lll@{}}
\toprule
Утверждение & Смысл & Где доказано \\
\midrule
$(3) \iff 0 \in \G$ & геометрический смысл условия & шаг 1 \\
$(3) \Rightarrow d_k = 0$ & достаточность & шаг 2 \\
$d_k = 0 \Rightarrow (3)$ & необходимость & шаг 3 \\
$0 \in \partial\varphi(0) \iff (3)$ & то же через условие оптимальности & шаг 4 \\
$(3) \not\Rightarrow \nabla f(\theta_k) = 0$ & ответ на вопрос задания & шаг 5 \\
\bottomrule
\end{tabular}
\end{center}

\section*{Шаг 1. Условие (3) эквивалентно $0 \in \G$}

Условие (3) выполнено ровно тогда, когда нулевой вектор лежит в множестве
неопределённости $\G$.

По определению $0 \in \G$ означает, что существует вектор $u$ с $\norm{u} \le \rho$
и $g + \Shalf u = 0$. Матрица $\Shalf$ обратима, поэтому это уравнение имеет
единственное решение $u = -\Smhalf g$. Остаётся проверить его норму:
\begin{equation*}
\norm{u}^2 = \big(\Smhalf g\big)^T\big(\Smhalf g\big)
= g^T\big(\Smhalf\big)^T\Smhalf g = g^T \Smhalf\Smhalf g = g^TS^{-1}g .
\end{equation*}
Следовательно, ограничение $\norm{u} \le \rho$ совпадает с (3):
\begin{equation*}
g^TS^{-1}g \le \rho^2 \iff 0 \in \G .
\end{equation*}

Величина $\sqrt{g^TS^{-1}g} = \norm{\Smhalf g}$ — это норма Махаланобиса вектора
$g$ относительно $S$. Она измеряет усреднённый градиент в единицах его собственной
оценённой неопределённости, то есть является отношением сигнал/шум.

\subsection*{Почему $\Smhalf\Smhalf = S^{-1}$}

По определению $\Shalf$ — симметричная положительно определённая матрица, квадрат
которой равен $S$, а $\Smhalf$ — обратная к ней. Берём обратную матрицу от обеих
частей равенства $S = \Shalf\Shalf$ и применяем правило $(AB)^{-1} = B^{-1}A^{-1}$:
\begin{equation*}
S^{-1} = \big(\Shalf\Shalf\big)^{-1} = \big(\Shalf\big)^{-1}\big(\Shalf\big)^{-1}
= \Smhalf\Smhalf .
\end{equation*}

Существование корня следует из спектрального разложения. Симметричная матрица
$S \succ 0$ записывается как $S = Q\Lambda Q^T$ с ортогональной $Q$ и диагональной
$\Lambda = \diag(\lambda_1,\dots,\lambda_p)$, $\lambda_i > 0$, и тогда
\begin{gather*}
\Shalf = Q\Lambda^{1/2}Q^T, \qquad \Smhalf = Q\Lambda^{-1/2}Q^T, \\
\Smhalf\Smhalf = Q\Lambda^{-1/2}\underbrace{Q^TQ}_{I}\Lambda^{-1/2}Q^T
= Q\Lambda^{-1}Q^T = S^{-1}.
\end{gather*}
Из этой же формулы видно, что $\Smhalf$ симметрична. Именно это использовано во
втором равенстве для $\norm{u}^2$ выше.

Для диагональной матрицы $S_k = \diag(s_{k,1},\dots,s_{k,p})$ из раздела 2 задания
условие (3) принимает покомпонентный вид
\begin{equation*}
\sum_{j=1}^{p} \frac{g_{k,j}^2}{s_{k,j}} \le \rho^2 .
\end{equation*}

\section*{Шаг 2. Достаточность: $(3) \Rightarrow d_k = 0$}

Если (3) выполнено, то $\varphi(d) > \varphi(0) = 0$ для всех $d \ne 0$, поэтому
$d = 0$ — единственная точка минимума.

Возьмём произвольное $d \neq 0$. Вставим $\Smhalf\Shalf = I$ между $g^T$ и $d$ и
применим неравенство Коши — Буняковского:
\begin{equation*}
g^Td = \big(\Smhalf g\big)^T\big(\Shalf d\big)
\;\ge\; -\norm{\Smhalf g}\,\norm{\Shalf d}
= -\sqrt{g^TS^{-1}g}\,\sqrt{d^TSd} \;\ge\; -\rho\sqrt{d^TSd}.
\end{equation*}
Последнее неравенство — это условие (3) после извлечения корня:
$\sqrt{g^TS^{-1}g} \le \rho$. Значит, первые два слагаемых $\varphi$ в сумме
неотрицательны, и остаётся только квадратичный штраф:
\begin{equation*}
\varphi(d) = \underbrace{g^Td + \rho\sqrt{d^TSd}}_{\ge\, 0} + \frac{\norm{d}^2}{2\alpha}
\;\ge\; \frac{\norm{d}^2}{2\alpha} \;>\; 0 = \varphi(0)
\qquad \text{для любого } d \ne 0 .
\end{equation*}
Строгое неравенство даёт не только минимальность нуля, но и единственность:
$d_k = 0$.

\subsection*{То же из минимаксной формы (1)}

Этот вывод не требует пункта (b). По шагу 1 условие (3) означает $0 \in \G$,
поэтому во внутреннем максимуме можно выбрать $h = 0$:
\begin{equation*}
\varphi(d) = \max_{h\in\G}\Big\{h^Td + \frac{\norm{d}^2}{2\alpha}\Big\}
\;\ge\; 0^Td + \frac{\norm{d}^2}{2\alpha} \;>\; 0 \qquad \text{при } d\ne 0 .
\end{equation*}
При $d = 0$ выражение под максимумом равно $0$ при любом $h$, поэтому
$\varphi(0) = 0$, и ноль снова единственная точка минимума.

Смысл: для каждого направления $d \ne 0$ в $\G$ есть допустимый градиент $h$ с
$h^Td \ge 0$. Ни одно направление не является гарантированным направлением
убывания, и любой шаг лишь платит штраф $\norm{d}^2/(2\alpha)$.

\section*{Шаг 3. Необходимость: $d_k = 0 \Rightarrow (3)$}

Если $d_k = 0$, то $\varphi$ неотрицательна всюду, и подстановка одного
специального направления даёт (3).

Пусть $d_k = 0$. Тогда $0$ — точка минимума $\varphi$, то есть
$\varphi(d) \ge \varphi(0) = 0$ для всех $d$. Обозначим
\begin{equation*}
\kappa = \sqrt{g^TS^{-1}g} \ge 0, \qquad c = \frac{\norm{S^{-1}g}^2}{2\alpha} \ge 0
\end{equation*}
и подставим направление $d(t) = -t\,S^{-1}g$ при $t > 0$. Три слагаемых $\varphi$
вычисляются так:
\begin{align*}
g^Td(t) &= -t\,g^TS^{-1}g = -t\kappa^2,\\
d(t)^TS\,d(t) &= t^2\,g^TS^{-1}SS^{-1}g = t^2\kappa^2
\;\;\Longrightarrow\;\; \rho\sqrt{d(t)^TS\,d(t)} = \rho t\kappa,\\
\frac{\norm{d(t)}^2}{2\alpha} &= t^2\,\frac{\norm{S^{-1}g}^2}{2\alpha} = t^2c.
\end{align*}
Неотрицательность $\varphi$ в точке $d(t)$ даёт
\begin{equation*}
\varphi(d(t)) = -t\kappa^2 + \rho t\kappa + t^2c \;\ge\; 0
\qquad \text{для всех } t > 0 .
\end{equation*}
Разделим на $t > 0$ и перейдём к пределу $t \to 0^+$:
\begin{equation*}
-\kappa^2 + \rho\kappa + tc \ge 0
\quad\xrightarrow[\;t\to 0^+\;]{}\quad \kappa(\rho - \kappa) \ge 0 .
\end{equation*}
Далее два случая:
\begin{itemize}
\item $\kappa = 0$: неравенство $\kappa \le \rho$ выполнено автоматически, так как
$\rho \ge 0$;
\item $\kappa > 0$: делим на $\kappa$ и получаем $\rho - \kappa \ge 0$, то есть
$\kappa \le \rho$.
\end{itemize}
В обоих случаях $\kappa \le \rho$. Обе части неотрицательны, поэтому возведение в
квадрат даёт \mbox{$g^TS^{-1}g \le \rho^2$}, то есть (3). \hfill $\blacksquare$

\subsection*{Эквивалентная форма: контрапозиция}

Импликации $A \Rightarrow B$ и $\neg B \Rightarrow \neg A$ логически эквивалентны:
обе ложны только в случае, когда $A$ истинно, а $B$ ложно. Здесь $A$ — это
«$d_k = 0$», $B$ — это (3). Поэтому утверждение шага 3 можно доказывать и в виде
«если (3) нарушено, то $d_k \ne 0$».

В этой форме вычисление то же. Пусть $\kappa > \rho$; тогда $g \ne 0$ и $c > 0$,
и для достаточно малых $t$
\begin{equation*}
\varphi(d(t)) = -t\,\kappa(\kappa-\rho) + t^2c \;<\; 0 = \varphi(0)
\qquad \text{при } 0 < t < \frac{\kappa(\kappa-\rho)}{c}
= \frac{2\alpha\,\kappa(\kappa-\rho)}{\norm{S^{-1}g}^2}.
\end{equation*}
Найдена точка, в которой $\varphi$ меньше, чем в нуле, поэтому ноль не является
точкой минимума и \mbox{$d_k \ne 0$}. Контрапозиция удобна тем, что для вывода
$d_k \ne 0$ достаточно предъявить одну такую точку.

Шаги 2 и 3 вместе дают эквивалентность
\begin{equation*}
d_k = 0 \iff g_k^TS_k^{-1}g_k \le \rho^2 .
\end{equation*}

\section*{Шаг 4. Тот же результат из условия оптимальности}

Условие (3) — это условие оптимальности $0 \in \partial\varphi(0)$, записанное в
точке излома $d = 0$; пункт (d) задания описывает оставшийся гладкий случай
$d_k \ne 0$.

Слагаемое $\rho\norm{\Shalf d}$ не дифференцируемо в точке $d = 0$, поэтому
градиентное уравнение из (d) там неприменимо и нужен субдифференциал. Функция
$d \mapsto \norm{\Shalf d}$ есть опорная функция эллипсоида
$\{\Shalf u : \norm{u} \le 1\}$:
\begin{equation*}
\norm{\Shalf d} = \max_{\norm{u}\le 1} u^T\Shalf d
= \max_{\norm{u}\le 1} \big(\Shalf u\big)^Td .
\end{equation*}
Субдифференциал опорной функции выпуклого компакта в нуле равен самому этому
множеству. Линейное слагаемое даёт $g$, а градиент квадратичного слагаемого в нуле
равен нулю, поэтому
\begin{equation*}
\partial\varphi(0) = \big\{\,g + \rho\,\Shalf u : \norm{u} \le 1\,\big\} = \G .
\end{equation*}
Функция $\varphi$ выпукла, поэтому по правилу Ферма включение
$0 \in \partial\varphi(0)$ необходимо и достаточно для того, чтобы $d = 0$ было
точкой минимума. Это включение читается как $0 \in \G$, что по шагу 1 равносильно
(3). Так шаги 2 и 3 получаются одной строкой.

Два следствия:
\begin{itemize}
\item \textbf{Порог не зависит от $\alpha_k$.} Шаг обучения масштабирует ненулевой
шаг, но не может превратить нулевой шаг в ненулевой и наоборот.
\item \textbf{Согласованность с (d).} При $S = sI$ условие (3) принимает вид
$\norm{g}^2/s \le \rho^2$, то есть $\rho\sqrt{s}/\norm{g} \ge 1$. Это в точности
случай, когда множитель в формуле из (d) обнуляется:
\begin{equation*}
d_k = -\alpha_k\Big(1 - \frac{\rho\sqrt{s}}{\norm{g_k}}\Big)_{+}\, g_k = 0
\iff \rho\sqrt{s} \ge \norm{g_k} .
\end{equation*}
\end{itemize}

\section*{Шаг 5. Следует ли из (3), что $\nabla f(\theta_k) = 0$}

Нет. Условие (3) говорит только, что нулевой градиент не исключён данными;
истинный градиент при этом может быть ненулевым.

\subsection*{Причина 1: в (3) входят только выборочные величины}

Условие (3) — утверждение о $g_k$, $S_k$ и настраиваемом параметре $\rho$.
Истинный градиент $\nabla f(\theta_k)$ в нём не участвует. Его можно сделать верным
в любой точке тремя способами: увеличить $\rho$, взять более шумный батч (больше
$S_k$) или уменьшить размер батча, так как $S_k$ оценивает ковариацию среднего по
батчу и убывает как $1/b$.

\subsection*{Причина 2: даже при $\nabla f(\theta_k) \in \G$ множество $\G$
содержит много векторов}

Предположим благоприятный случай из пункта (e): $\nabla f(\theta_k) \in \G$.
Условие (3) добавляет к этому $0 \in \G$. При $\rho > 0$ множество $\G$ —
полноразмерный эллипсоид, и вместе с нулём в нём лежит бесконечно много ненулевых
векторов; $\nabla f(\theta_k)$ может быть любым из них.

Правильное прочтение — статистическое. Условие (3) означает, что гипотеза
$\nabla f(\theta_k) = 0$ «не отвергается» при радиусе $\rho$, но не что она
доказана.

\subsection*{Причина 3: контрпример, построенный по формуле (4)}

Возьмём $f(\theta) = \tfrac12\theta^2$ в точке $\theta_k = 1$. Тогда
$\nabla f(\theta_k) = 1 \ne 0$, а минимум находится в $\theta^\star = 0$. Пусть батч
состоит из $b = 2$ зашумлённых наблюдений градиента: $v_{k,1} = 3$ и
$v_{k,2} = -1$. Все величины строятся по формулам задания:
\begin{align*}
g_k &= \tfrac12\big(3 + (-1)\big) = 1,\\
\widehat s_k &= \frac{1}{b(b-1)}\sum_i (v_{k,i} - g_k)^2
= \frac{1}{2}\big(2^2 + (-2)^2\big) = 4,\\
\varepsilon_k &= 10^{-3}\max\{4,\,10^{-8}\} = 0.004,\\
S_k &= s_k = \widehat s_k + \varepsilon_k = 4.004 .
\end{align*}
Проверим (3) при умеренном радиусе $\rho = 1$:
\begin{equation*}
g_k^TS_k^{-1}g_k = \frac{1^2}{4.004} \approx 0.2498 \;\le\; 1 = \rho^2 .
\end{equation*}
Условие (3) выполнено, поэтому $d_k = 0$ и $\theta_{k+1} = \theta_k = 1$. При этом:
\begin{itemize}
\item $\nabla f(\theta_k) = 1 \ne 0$, точка не стационарна;
\item оценка градиента даже точна, $g_k = \nabla f(\theta_k)$, и включение
$\nabla f(\theta_k) \in \G$ выполнено;
\item обычный SGD сделал бы правильный шаг $-\alpha_k$ в сторону минимума.
\end{itemize}

Стандартная ошибка среднего здесь $\sqrt{s_k} \approx 2$, а само среднее равно $1$,
то есть отстоит от нуля на половину стандартной ошибки. По двум таким наблюдениям
нельзя отличить истинный градиент $1$ от истинного градиента $0$, и R-SGD
отказывается делать шаг.

Проверка по формуле из (d): $\rho\sqrt{s_k}/|g_k| \approx 2.001 > 1$, множитель
$(1 - 2.001)_+$ равен $0$. Численная минимизация $\varphi$ для этих данных также
возвращает $d = 0$ при $\rho = 1$ и ненулевой шаг при $\rho = 0.4$, где (3)
нарушено.

\subsection*{Когда импликация всё же верна}

Единственный случай — $\rho = 0$ вместе с включением $\nabla f(\theta_k) \in \G$.
Тогда (3) даёт $g_k = 0$, множество $\G$ сжимается в $\{0\}$, и отсюда
$\nabla f(\theta_k) = 0$. Без включения даже $g_k = 0$ ничего не доказывает, потому
что $g_k$ — зашумлённая оценка по мини-батчу.

Обратная импликация тоже неверна: из $\nabla f(\theta_k) = 0$ не следует (3).
Зашумлённый $g_k$ может выйти за пределы шара Махаланобиса радиуса $\rho$ и дать
ненулевой шаг в точной стационарной точке.

\section*{Интерпретация и проверка}

Нулевой шаг возникает там, где градиент мал по сравнению с его оценённым шумом, а
это происходит раньше, чем достигается стационарная точка.

\begin{itemize}
\item \textbf{Нулевые шаги до минимума.} Вблизи минимума $\norm{\nabla f}$
уменьшается, а шум градиента нет, поэтому $g^TS^{-1}g$ опускается ниже $\rho^2$ до
достижения стационарности. Для квадратичного примера из задания,
$f(\theta) = \tfrac12(\theta_1^2 + 4\theta_2^2)$, при точных градиентах и $S = sI$
область нулевого шага — эллипс $\theta_1^2 + 16\theta_2^2 \le \rho^2 s$ вокруг
$\theta^\star = 0$. Это материал для раздела 5(b).
\item \textbf{Нулевой шаг не обязательно навсегда.} В стохастической постановке на
следующей итерации берётся новый батч, $g_{k+1}$ и $S_{k+1}$ меняются, и (3) может
нарушиться в той же точке $\theta$. При фиксированных $g$ и $S$ метод остановился
бы окончательно.
\item \textbf{Роль размера батча.} Оценка (4) убывает примерно как $1/b$, поэтому
при том же $g_k$ величина $g^TS^{-1}g$ растёт примерно линейно по $b$. Большой батч
сужает область нулевого шага, маленький расширяет.
\item \textbf{Роль $\rho$ и $\alpha_k$.} Частота нулевых шагов определяется парой
$(\rho, S_k)$ и не зависит от $\alpha_k$ напрямую (шаг 4).
\end{itemize}

\subsection*{Численная проверка}

Эквивалентность $d_k = 0 \iff (3)$ подтверждена на 300 случайных задачах без
единого расхождения. В каждой задаче размерность $p$ от 1 до 3, матрица $S$
диагональная, а радиус $\rho$ брался равным $\kappa$, умноженному на $0.5$, $0.9$,
$0.99$, $1.01$, $1.1$ или $2$, чтобы покрыть обе стороны порога и его окрестность.
Функция $\varphi$ минимизировалась методом Нелдера — Мида из 20 случайных
начальных точек.

Для контрпримера из шага 5 ($g_k = 1$, $s_k = 4.004$, $\alpha_k = 0.1$) численный
минимум такой:

\begin{center}
\begin{tabular}{@{}llll@{}}
\toprule
$\rho$ & $g^TS^{-1}g$ против $\rho^2$ & Численный $d_k$ & Формула из (d) \\
\midrule
$1$ & $0.2498 \le 1$, (3) выполнено & $0$ & $0$ \\
$0.4$ & $0.2498 > 0.16$, (3) нарушено & $-0.01996$ & $-0.01996$ \\
\bottomrule
\end{tabular}
\end{center}

Для раздела 2(b) та же проверка делается решателем задачи (2): он должен возвращать
$d = 0$ с точностью до допуска ровно тогда, когда $g^TS^{-1}g \le \rho^2$, включая
значения $\rho$ чуть выше и чуть ниже $\kappa$.

\end{document}
